\section{Arquitectura de Datos: Proceso de Extracción, Transformación y Carga (ETL)}

Para procesar los datos, usamos un modelo de Esquema de Constelaci\'on,
y luego para el an\'alisis, juntamos todos los datos para la tabla del
an\'alisis correspondiente. 
La idea es tener una buena base para poder analizar con python 
de manera eficiente. Las tablas estandarizadas est\'an separadas en
tablas de dimensi\'on (\texttt{dim\_*.csv}) y tablas de hechos 
(\texttt{fact\_*.csv}).

\subsection{Carga de datos}

Los datasets usados son los siguientes:
\begin{itemize}[noitemsep]
    \item \href{https://catalogodatos.gub.uy/dataset/velocidad-promedio-vehicular-en-las-principales-avenidas-de-montevideo}{Velocidad promedio}
    \item \href{https://catalogodatos.gub.uy/dataset/conteo-vehicular-en-las-principales-avenidas-de-montevideo}{Conteo vehicular}
    \item \href{https://catalogodatos.gub.uy/dataset/base-anual-de-personas-lesionadas-en-siniestros-de-transito}{Lesionados en siniestros}
\end{itemize}

La carga de datos se realiza autom\'aticamente con el script \texttt{descargar\_con\_fechas.py}.

\subsection{Construcción de Dimensiones (\texttt{build\_dim\_*.py})}

Los módulos de dimensión procesan los datos crudos para extraer entidades
descriptivas, normalizar sus atributos y asignarles claves 
(IDs numéricos enteros). 
Este proceso garantiza la integridad referencial del modelo de datos y elimina
la redundancia de datos almacenados.

\subsubsection*{Script: \texttt{build\_dim\_calle.py}}
Este script genera un catálogo único y estandarizado de todas las calles de la
ciudad. 
Toma como entrada la dimensión de ubicaciones base y extrae la columna
descriptiva original, la cual frecuentemente contiene cadenas compuestas de 
intersecciones viales (ej. ``Avenida A / Calle B'' o ``Calle C esq. Calle D'').

El procesamiento ejecuta las siguientes limpiezas:
\begin{itemize}[noitemsep]
    \item \textbf{Fragmentación de Cadenas}: 
        Divide el texto iterando sobre delimitadores como `` / '' y `` esq. ''.
    \item \textbf{Estandarización de Texto}:
        Elimina espacios en blanco residuales y lleva todo el texto a letras 
        mayúsculas.
    \item \textbf{Filtrado de Anomalías}:
        Aplica una expresión regular 
        (\textrm{r`\string^-?\textbackslash d+\textbackslash.?\textbackslash d*\$'}) 
        para detectar y descartar fragmentos que consisten puramente en números,
        los cuales representan errores de tipeo o marcadores kilométricos 
        inválidos como nombres de calles.
    \item \textbf{Indexación}: 
        Agrupa los valores para asegurar la unicidad y asigna un identificador
        secuencial (\texttt{calle\_id}), exportando el resultado al archivo
        \texttt{dim\_calle.csv}.
\end{itemize}

\subsubsection*{Script: \texttt{build\_dim\_fecha.py}}

Este script crea una tabla de calendario completa para desglosar cualquier fecha
por año, mes o día. 
No asume un solo formato de datos porque los archivos del gobierno siempre
varían en ese sentido.
\begin{itemize}[noitemsep]
    \item \textbf{Consolidación Multi-Fuente}:
        Itera sobre los directorios de velocidad, conteo vehicular y siniestros.
        Para los siniestros, aplica un mapeo de columnas específico 
        (SIN\_MAPPING) debido a discrepancias en el esquema de origen.
    \item \textbf{Conversión Temporal}:
        Convierte el texto crudo al tipo datetime.
    \item \textbf{Derivación de Atributos}:
        Tras obtener un listado único y cronológico de todas las fechas 
        registradas, el algoritmo extrae directamente del objeto datetime tres 
        nuevas columnas de análisis: el año (anio), el mes numérico (mes) y el 
        día (dia). Finalmente, asigna una clave primaria.
\end{itemize}

\subsubsection*{Script: \texttt{build\_dim\_detector.py}}
\begin{itemize}[noitemsep]
    \item \textbf{Extracción de Metadata}:
        Selecciona el código natural del detector y sus descriptores geográficos
        (avenida, intersección anterior e intersección siguiente).
    \item \textbf{Casteo de Precisión}:
        Fuerza la conversión de los tipos de datos de latitud y longitud a
        Float64 explícitamente.
    \item \textbf{Renombramiento Estructural}:
        Limpia los nombres de las columnas, eliminando prefijos innecesarios 
        (\texttt{dsc\_}), y exporta la tabla única de sensores con su 
        respectivo ID.
\end{itemize}

\subsubsection*{Script: \texttt{build\_dim\_ubicacion.py}}

Este script unifica dos tipos de coordenadas a uno, consolidando los datos en
una tabla diferenciada por el \texttt{tipo}.
\begin{itemize}[noitemsep]
    \item \textbf{Procesamiento de Coordenadas GPS}: 
        Para los datos de flujo vehicular, extrae directamente la latitud y
        longitud y asigna el identificador categórico \texttt{tipo = `avenida'}.
    \item \textbf{Transformación de Coordenadas UTM}:
        El dataset de siniestros no provee latitud y longitud, sino coordenadas
        planas en formato UTM (x, y). 
        El script usa la funci\'on \texttt{utm\_to\_latlon} para pasar de UTM a 
        latitud y longitud. 
        A estos puntos les asigna el indicador \texttt{tipo = `calle\_siniestro'}.
    \item \textbf{Agrupación y Fusión de Descripciones}:
        Consolida todos los puntos en un único DataFrame.
        Si múltiples sensores o accidentes comparten las mismas coordenadas, 
        agrupa el registro fusionando las descripciones de texto para no 
        perder contexto, asignando luego la llave primaria final.
\end{itemize}

\subsection{Construcción de Hechos (\texttt{build\_fact\_*.py})}

Los módulos de hechos se encargan de consolidar el volumen alto de datos del
sistema. Su función principal es reemplazar los datos que ya consolidamos en 
las tablas de dimensi\'on por los identificadores correspondientes.

\subsubsection*{Script: \texttt{build\_fact\_medicion.py}}
\begin{itemize}[noitemsep]
    \item \textbf{Cruce de Orígenes}:
        En cada iteración mensual, normaliza los textos de fecha y extrae la
        hora y el minuto como números enteros. Utiliza una llave compuesta de 5
        variables (\texttt{[`fecha', `hora', `minuto', `cod\_detector', `id\_carril']}) 
        para ejecutar un Left Join, la velocidad promedio de un vehículo con
        el volumen de la calle en ese minuto exacto.
    \item \textbf{Reconstrucción Temporal}:
        Concatena las variables de hora y minuto que estaban separadas en un
        formato de texto estructurado (\texttt{"HH:MM:00"}) y lo convierte en
        un tipo de dato nativo de tiempo (\texttt{hora\_t}).
    \item \textbf{Mapeo de Claves y Streaming}:
        Cruza el resultado con las dimensiones en memoria para obtener
        el \texttt{fecha\_id} y \texttt{detector\_id}. Elimina todas las 
        columnas descriptivas dejando solo números. 
\end{itemize}

\subsubsection*{Script: \texttt{build\_fact\_siniestros.py}}

Se encarga de modelar la tabla de transacciones de accidentes de tránsito, 
conectando la metadata categórica del siniestro con las dimensiones espacial
y temporal.
\begin{itemize}[noitemsep]
    \item \textbf{Estandarización de Esquema y Texto}: 
        Primero, se cargan los datos y se utiliza el diccionario SIN\_MAPPING
        para renombrar las columnas y ajustarlas a la convención definida en el
        proyecto. 
        Adicionalmente, se aplica una limpieza (\textit{strip}) a 
        las columnas de texto para eliminar cualquier espacio en blanco
        innecesario al inicio o al final de las palabras.
    \item \textbf{Geoprocesamiento Integrado}:
        Al igual que la dimensión de ubicación, este script de hechos repite la
        transformación geométrica \texttt{utm\_to\_latlon}.
    \item \textbf{Cruces}:
        Ejecuta un cruce conectando la fecha del siniestro con
        \texttt{dim\_fecha}, y un cruce espacial sobre la latitud y longitud 
        para recuperar el \texttt{ubicacion\_id} generado en la etapa de 
        dimensiones.
    \item \textbf{Limpieza Final y Exportación}:
        Desecha las coordenadas crudas y las fechas en formato texto, 
        preservando únicamente las claves (\textit{IDs}) y las
        variables del accidente (si el conductor 
        utilizaba casco o cinturón, sexo, edad, tipo de siniestro). 
\end{itemize}
